\documentclass[11pt,a4paper]{article}

\usepackage[utf8]{inputenc}
\usepackage[T1]{fontenc}
\usepackage{lmodern}
\usepackage{amsmath,amssymb,bm}
\usepackage{enumitem}
\usepackage{hyperref}
\usepackage{verbatim}
\usepackage{microtype}
\usepackage[margin=2.5cm]{geometry}
\usepackage{xcolor}
\usepackage{booktabs}
\usepackage{array}
\newcolumntype{L}[1]{>{\raggedright\arraybackslash}p{#1}}

\setlength{\parskip}{0.4em}
\setlength{\parindent}{0pt}
\renewcommand{\labelenumi}{(\alph{enumi})}
\definecolor{panel}{gray}{0.94}
\newcommand{\R}{\mathbb{R}}
\newcommand{\vect}[1]{\bm{#1}}

\newsavebox{\ctxboxsv}
\newenvironment{ctxbox}[1]{%
  \par\smallskip\noindent
  \begin{lrbox}{\ctxboxsv}
  \begin{minipage}{\dimexpr\linewidth-2\fboxsep-2\fboxrule\relax}
  \textbf{#1}\par\smallskip
}{%
  \end{minipage}
  \end{lrbox}%
  \noindent\colorbox{panel}{\usebox{\ctxboxsv}}\par\smallskip
}

\title{\bfseries Lab Assignment 2\\[4pt]
\Large Three-Dimensional Rotations and Their Representations}
\author{Maths II -- Bachelor's Degree in Video Game Design and Development\\CITM - UPC}
\date{Autumn 2026}

\begin{document}
\maketitle

\begin{ctxbox}{Abstract}
This assignment extends the vector and matrix algebra developed in lab 1 to
three-dimensional rotations. Students implement conversions between rotation
matrices, axis-angle descriptions, intrinsic Euler ZYX angles, and quaternions.
They also rotate vectors and measure the arithmetic cost of selected operations.
The goal is to connect geometric definitions with reliable C++ implementations,
including representation ambiguities, singular cases, and input validation.
An interactive app and automated tests provide complementary validation.
\end{ctxbox}

\textbf{Repository:} \url{https://github.com/jvllahi/math_2_lab}

\section*{Scope and Conventions}

All work is in the \texttt{lab2/} project. Implement the TODO methods in
\texttt{assignment/Matrix3x3.cpp} and \texttt{assignment/Quat.cpp}; their headers
define the public contracts. The generic algebra and supplied helper functions
are not new student exercises.

\begin{itemize}[leftmargin=*]
    \item Rotations are active and right-handed, acting on column vectors.
    Matrix entries are stored in row-major order.
    \item Function angles are in radians; the app displays and accepts degrees.
    \item Euler angles follow intrinsic ZYX, with yaw about Z, pitch about the
    local Y axis, and roll about the local X axis.
    \item Quaternion components are ordered \texttt{(s,x,y,z)}: scalar first.
    \item Rotation checks use absolute tolerance $10^{-6}$ for orthogonality
    and determinant. The app accepts numerical output error up to $5\cdot10^{-6}$;
    automated tests also check stricter tolerances for selected cases.
\end{itemize}

The following equations describe the mathematical objects and conventions.
Deriving the conversion algorithms and handling their numerical edge cases are
part of the assignment; this handout does not provide solution code or pseudocode.

\section*{Building on Lab 1}

\textbf{Before configuring or building lab2, create a folder named
\texttt{algebra/} inside \texttt{lab2/}.} Copy your own completed
\texttt{LinAlgebra.hpp} and \texttt{LinAlgebra.cpp} from the first assignment
into it, keeping their filenames and original public API. This folder is
intentionally not distributed and is ignored by Git; do not use an instructor
reference solution.

Your copied files provide the generic \texttt{Vector}, \texttt{Matrix}, and
\texttt{Operations} needed by lab2. They are required for both app and test
builds, regardless of the implementation preset.
\texttt{Vec3} contains a three-element \texttt{Vector}; \texttt{Matrix3x3}
contains a $3\times3$ \texttt{Matrix}. Their fixed-size arithmetic operators
reuse the generic implementation.

The supplied \texttt{helpers/} include matrix allocation and row swaps, the
Euclidean norm, fixed-size construction and access, basic geometry, quaternion
normalization, and independent validation. Use vector indices $0,1,2$ and
\texttt{At(row,col)} for matrix entries. Default vectors and matrices are zero;
\texttt{Matrix3x3::Identity()} constructs the identity.
Do not copy duplicate implementations of the utilities already supplied in
\texttt{helpers/Numerics.cpp}.

\subsection*{Why use algebraic expressions?}

Before writing component-by-component loops, consider whether the matrix and
vector operators from lab 1 can express the algebra you need. They are also
available for \texttt{Matrix3x3} and \texttt{Vec3}, with the same natural syntax.

Reusing these operations can keep your code closer to the mathematical notation,
reduce repetitive arithmetic, and make it easier to read and check. Choosing how
to apply them to the rotation methods is part of the assignment.

\section*{Exercise 1: Rotation Matrices and Axis-Angle}

\subsection*{a) Mathematical Summary}

A proper rotation matrix belongs to
\[
SO(3)=\{R\in\R^{3\times3}:R^{\mathsf T}R=I,\ \det R=1\}.
\]
Its columns form a right-handed orthonormal basis. A rotated vector is
$\vect{v}'=R\vect{v}$, and rotation preserves lengths and dot products.
Orthogonality alone also permits reflections, which have determinant $-1$.

An axis-angle description consists of a unit axis $\vect{u}$ and a physical
rotation angle $\theta$. Rodrigues' geometric relation is
\[
R\vect{v}
=\vect{v}\cos\theta
+(\vect{u}\times\vect{v})\sin\theta
+\vect{u}(\vect{u}\cdot\vect{v})(1-\cos\theta).
\]
The axis is unchanged by the rotation. Equivalent descriptions are not unique:
angles can differ by full turns, and reversing both axis and angle describes
the same rotation. For output, use a unit axis and $\theta\in[0,\pi]$.
The identity has no unique axis; at a half-turn, opposite axes are equivalent.
These ambiguities must not be confused with incorrect rotations.

\subsection*{b) Implementation}

Implement in \texttt{assignment/Matrix3x3.cpp}:
\begin{itemize}[leftmargin=*]
    \item \texttt{IsRotation}: return whether the matrix is a proper rotation
    within the specified tolerance. Non-finite entries return \texttt{false}.
    \item \texttt{RotationAxisAngle}: construct a rotation from an axis and
    an angle. Normalize the finite, nonzero axis; reject a zero axis even when
    the angle is zero.
    \item \texttt{ToAxisAngle}: extract a unit axis and an angle in $[0,\pi]$.
    Handle identity, half-turns, and nearby angles reliably.
    \item \texttt{Rotate}: rotate a finite vector using a valid rotation
    matrix and accumulate the scalar arithmetic counts of the product kernel.
\end{itemize}

Except for the boolean \texttt{IsRotation} check, invalid matrix, axis, vector,
or angle inputs must produce \texttt{std::invalid\_argument}. Choose a
deterministic unit axis for identity output and explain the choice.

\subsection*{c) Report}
\begin{enumerate}[leftmargin=*]
    \item Why must a proper rotation satisfy both orthogonality and determinant
    $+1$? Provide a matrix that satisfies one condition but is not a rotation.
    \item Explain geometrically why rotation preserves vector length and leaves
    its axis unchanged. Show numerical evidence from one test.
    \item Why are identity and half-turn rotations special for axis-angle
    extraction? Explain the ambiguities and how your implementation handles them.
    \item Report one normal and one near-half-turn round trip. State the input,
    the returned parameters, and the reconstructed matrix error.
\end{enumerate}

\section*{Exercise 2: Intrinsic Euler ZYX Angles}

\subsection*{a) Mathematical Summary}

Euler angles describe orientation using three elementary rotations. This
assignment fixes the convention
\[
R(\psi,\theta,\phi)=R_z(\psi)R_y(\theta)R_x(\phi),
\]
where $\psi$ is yaw, $\theta$ is pitch, and $\phi$ is roll. The matrices
$R_x$, $R_y$, and $R_z$ are right-handed rotations about the corresponding
coordinate axes. Matrix multiplication is not generally commutative, so the
axis order and the active column-vector convention matter.

Different Euler triples can represent the same orientation. Canonical outputs
must satisfy
\[
\psi,\phi\in[-\pi,\pi],\qquad
\theta\in[-\pi/2,\pi/2].
\]
At $\theta=\pm\pi/2$, two rotational degrees of freedom become coupled:
this is gimbal lock. A unique yaw/roll pair cannot be recovered. For a
reproducible output convention in this assignment, set roll to zero at exact
gimbal lock and return an equivalent orientation.

\subsection*{b) Implementation}

Implement in \texttt{assignment/Matrix3x3.cpp}:
\begin{itemize}[leftmargin=*]
    \item \texttt{FromEulerZYX}: construct the specified orientation from
    finite yaw, pitch, and roll inputs in radians.
    \item \texttt{ToEulerZYX}: extract canonical Euler angles from a valid
    rotation matrix, including both signs of gimbal lock and nearby cases.
\end{itemize}

Reject non-finite input angles or non-rotation matrices with
\texttt{std::invalid\_argument}. Validate equivalent orientations rather than
requiring an extracted Euler triple to equal a noncanonical input triple.

\subsection*{c) Report}
\begin{enumerate}[leftmargin=*]
    \item Explain the chosen intrinsic ZYX convention. Why would changing the
    multiplication order generally change the orientation?
    \item What is lost at gimbal lock? Explain why different yaw/roll pairs can
    describe the same orientation, and state your output convention.
    \item Compare a normal round trip with pitch at $+90^\circ$, $-90^\circ$,
    and close to either value. Include reconstructed matrix errors.
    \item Give two distinct Euler triples representing the same orientation.
    Explain why comparing angles directly is not a sufficient correctness test.
\end{enumerate}

\section*{Exercise 3: Quaternions and Rotation}

\subsection*{a) Mathematical Summary}

A quaternion has a scalar part and a three-dimensional vector part:
\[
q=s+x\,i+y\,j+z\,k=(s,\vect{u}),\qquad
\vect{u}=(x,y,z).
\]
The units satisfy $i^2=j^2=k^2=ijk=-1$. For
$p=(s,\vect{u})$ and $q=(t,\vect{w})$, the Hamilton product is defined by
\[
pq=(st-\vect{u}\cdot\vect{w},\
s\vect{w}+t\vect{u}+\vect{u}\times\vect{w}).
\]
It is associative but not generally commutative.

The norm and conjugate are
\[
\|q\|=\sqrt{s^2+x^2+y^2+z^2},\qquad
q^*=(s,-x,-y,-z).
\]
Unit quaternions encode rotations. The axis-angle representation corresponds to
\[
q=\left(\cos\frac{\theta}{2},\
\vect{u}\sin\frac{\theta}{2}\right),
\qquad \|\vect{u}\|=1.
\]
For a unit quaternion, the vector part of $q(0,\vect{v})q^*$ is the rotated
vector. The quaternions $q$ and $-q$ represent the same rotation.
The product $q_1q_2$ composes rotations with $q_2$ acting first.

\subsection*{b) Implementation}

Implement in \texttt{assignment/Quat.cpp}:
\begin{itemize}[leftmargin=*]
    \item \texttt{Multiply}: compute the Hamilton product without normalizing
    either operand, and accumulate scalar arithmetic counts. Finite zero
    quaternions are valid operands for this algebraic operation.
    \item \texttt{Rotate}: normalize the quaternion, rotate a finite vector,
    and accumulate only the rotation kernel's scalar arithmetic counts.
    \item \texttt{FromAxisAngle} and \texttt{ToAxisAngle}: convert between
    unit quaternions and axis-angle descriptions with the output conventions
    of exercise 1.
    \item \texttt{FromMatrix3x3} and \texttt{ToMatrix3x3}: convert between
    unit quaternions and proper rotation matrices.
\end{itemize}

\texttt{Normalized} is supplied in the helpers. Rotation and conversion methods
must normalize finite, nonzero quaternion inputs; zero quaternions do not
describe rotations. Axis inputs must be finite and nonzero; angle inputs must
be finite. Invalid inputs throw \texttt{std::invalid\_argument}.
Matrix-to-quaternion and axis-to-quaternion outputs must have unit norm.
Axis-angle extraction must treat $q$ and $-q$ consistently.

\subsection*{c) Report}
\begin{enumerate}[leftmargin=*]
    \item Explain the role of the half-angle and why $q$ and $-q$ describe
    the same rotation. Demonstrate the latter numerically.
    \item Give a pair of quaternion rotations whose composition changes when
    the operands are swapped. Relate the result to matrix composition.
    \item Why is normalization appropriate for rotation/conversion methods,
    but not for an arbitrary Hamilton product? What happens with a zero quaternion?
    \item Compare rotating the same vector using a matrix and a quaternion.
    Report output error and each arithmetic count, and justify any difference
    from the app's advisory standard counts.
    \item Demonstrate a matrix--quaternion--matrix round trip for a normal
    rotation and a half-turn. Explain why quaternion components alone are not
    a sufficient equivalence check.
\end{enumerate}

\section*{Validation and Arithmetic Counts}

Use \texttt{Operations} to accumulate scalar additions, subtractions,
multiplications, and divisions in \texttt{Matrix3x3::Rotate},
\texttt{Quat::Multiply}, and \texttt{Quat::Rotate}. Do not reset a counter
provided by the caller.

Exclude index arithmetic, comparisons, control flow, input validation,
normalization, and independent reference calculations. The ordinary overloaded
algebra operators do not update these counters automatically. Account for the
arithmetic actually performed by the generic operations you call, including
additions to a zero accumulator where applicable.

The app's count comparison is advisory, separate from numerical correctness.
Explain your counting model and any alternative algorithm in the report.
Elapsed time is informational, not a standalone proof of correctness.

Use round trips and geometric invariants to compare orientations. A small
error in recovered parameters is not always meaningful near a non-unique
representation. Report the error in reconstructed matrices or rotated vectors,
and verify the required output ranges and unit norms separately.

\section*{How to Test the Implementation}

First copy your completed lab1 algebra into \texttt{lab2/algebra/}, then run
the supplied presets from the \texttt{lab2/} folder:
\begin{verbatim}
cmake --preset default-tests
cmake --build --preset default-tests --parallel
ctest --preset default-tests
\end{verbatim}

The test executable is \texttt{build-tests/bin/lab2\_tests}, with
\texttt{.exe} on Windows. Tests need no display. Unfinished rotation TODOs are
expected to fail; the supplied geometry tests should already pass.
Add focused tests for your own edge cases rather than relying only on the
provided examples.

Also build and run the interactive app:
\begin{verbatim}
cmake --preset default
cmake --build --preset default --parallel
./build/bin/lab2
\end{verbatim}
On Windows PowerShell, run \verb|.\build\bin\lab2.exe| instead.
Use \textbf{Operation}, \textbf{Example}, \textbf{Load example}, and
\textbf{Run operation}. Check status and details, numerical output, error,
and applicable counts. A rejected invalid input can correctly produce no
numerical output. See \texttt{README.md} for the full UI workflow.

\begin{center}
\begin{tabular}{L{0.30\linewidth}L{0.60\linewidth}}
\toprule
\textbf{Case} & \textbf{What to verify}\\
\midrule
Normal rotations & Known orientations, vector length, valid matrix, and round trips.\\
Identity & Equivalent identity orientation and valid, deterministic output parameters.\\
Half-turn and nearby angles & Stable axis-angle and quaternion conversions with unit outputs.\\
Gimbal lock and nearby pitch & Both signs, canonical Euler ranges, and equivalent reconstructed rotation.\\
Quaternion sign and scale & Equivalent rotation for opposite signs and nonzero scaled inputs.\\
Invalid inputs & Reflection/non-rotation checks; zero axis or quaternion, and non-finite values rejected where required.\\
\bottomrule
\end{tabular}
\end{center}

Record reproducible inputs, expected properties, actual results, and tolerances.
Include automated test results and at least one student-authored edge case for
each exercise in your report.

\section*{Evaluation Criteria}

The same weighting structure as lab 1 applies:
\begin{itemize}[leftmargin=*]
    \item Exercise 1 (35\%): Implementation 25\% + Report 10\%.
    \item Exercise 2 (35\%): Implementation 25\% + Report 10\%.
    \item Exercise 3 (30\%): Implementation 20\% + Report 10\%.
\end{itemize}

Evaluation considers mathematical correctness, adherence to the API contracts,
singular-case handling, validation, justified arithmetic counts, reproducible
tests, and the explanations requested in each report section.

The assignment is optional. If it is not submitted, the final lab grade is the
lab exam grade. If it is submitted, the final lab grade is the greater of the
exam grade and the average of the assignment and exam grades, so submission
does not penalize the student.

The code must work correctly; otherwise, the corresponding exercise receives
zero. If plagiarism is detected, the final grade is zero.

\section*{Academic Integrity and AI Use}

Using AI tools (LLMs, code assistants) as tutoring support is allowed.
Submitting generated code or report text as your own work without disclosure
is not allowed.

\begin{itemize}[leftmargin=*]
    \item If AI support was used, include \texttt{AI\_USAGE.md} in the submitted zip.
    \item List the tool name, what it was used for, and the final prompts that
    directly influenced the submitted code or report.
    \item You must be able to explain and justify every submitted function and test.
\end{itemize}

During review, students will be asked to explain selected code fragments or test
decisions. A submitted part that a student cannot explain will be considered
invalid for grading.

\section*{Submission}

Groups may contain up to three people. Submit through ATENEA using the naming
pattern below, omitting unused member slots for smaller groups:
\begin{center}
\small\texttt{Surname1Name1\_Surname2Name2\_Surname3Name3\_Lab2.zip}
\end{center}

Include:
\begin{itemize}[leftmargin=*]
    \item A buildable \texttt{lab2/} project preserving its relative folder
    structure: completed assignment sources and headers, your own previous
    implementation in \texttt{algebra/}, provided \texttt{helpers/} and
    \texttt{app/}, tests, and CMake files.
    Include any student-authored tests and necessary build changes.
    \item A PDF report of 2--4 pages, with a section for each exercise answering
    its report questions. Identify the group members and include numerical
    evidence, count comparisons, test results, and conclusions.
    \item \texttt{AI\_USAGE.md}, if applicable.
\end{itemize}

Do not include the provided \texttt{solution/} folder, build directories,
executables, or downloaded dependencies. Your submitted implementation must
build with the default assignment preset from a fresh project directory.
The submission is evaluated jointly with the group members; the lab exam is
individual.

\section*{Important Dates}
\noindent\fbox{%
\parbox{\dimexpr\linewidth-2\fboxsep-2\fboxrule\relax}{%
\centering\bfseries
Lab 2 Assignment Submission Deadline: To be announced.\\
Term Exam: To be announced.\\
Lab Practical Exam: To be announced.
}%
}

\end{document}
